\documentclass[twocolumn,11pt,a4paper]{article}
\usepackage[utf8]{inputenc}
\usepackage[T1]{fontenc}
\usepackage[margin=2.3cm]{geometry}
\usepackage{amsmath,amssymb,amsthm,mathtools}
\usepackage{bm}
\usepackage{booktabs}
\usepackage{array}
\usepackage{graphicx}
\usepackage{caption}
\captionsetup{font=small,labelfont=bf}
\usepackage{subcaption}
\usepackage{hyperref}
\usepackage{natbib}
\bibliographystyle{unsrtnat}
\usepackage{lmodern}
\usepackage[expansion=false]{microtype}
\usepackage{algorithm}
\usepackage{algpseudocode}
\usepackage{xcolor}

\hypersetup{colorlinks=true, linkcolor=blue!65!black, citecolor=blue!65!black, urlcolor=blue!65!black}

\newtheorem{theorem}{Theorem}[section]
\newtheorem{proposition}[theorem]{Proposition}
\newtheorem{lemma}[theorem]{Lemma}
\newtheorem{corollary}[theorem]{Corollary}
\theoremstyle{definition}
\newtheorem{definition}[theorem]{Definition}
\theoremstyle{remark}
\newtheorem{remark}[theorem]{Remark}

\newcommand{\bigO}{\mathcal{O}}
\newcommand{\R}{\mathbb{R}}

\title{\textbf{Spectral Pathogen Surveillance: Real-Time Identification,
Tropism Prediction, and Infection Dynamics Monitoring via
Oscillator Interference Analysis}}

\author{Kundai Farai Sachikonye\\
AIMe Registry for Artificial Intelligence\\
\texttt{kundai.sachikonye@bitspark.com}}

\date{}

\begin{document}

\maketitle

\begin{abstract}
Identifying pathogens, predicting cellular tropism, forecasting variant
emergence, monitoring antibiotic resistance, and tracking the host response to
infection are distinct problems currently solved by separate computational
pipelines, each requiring large databases, expensive computation, and
domain-specific software. We show that all five reduce to a single operation:
measuring constructive interference between the spectral representations of
biological sequences. Under a channelised discrete Fourier representation,
pathogen identity is a spectral fingerprint, tropism is the spectral overlap
between pathogen attachment proteins and host surface receptors, immune escape
variants are mutations that shift the pathogen spectrum out of antibody
recognition bands while retaining receptor-binding overlap, antibiotic
resistance is enzyme spectral evasion of drug spectra, and infection-driven
disease progression is the time-evolution of host oscillator coherence under
spectral perturbation by pathogen proteins. A unified GPU fragment-shader
implementation synthesises and compares spectra on demand with no pre-built
database and no stored structural coordinates. On synthetic benchmarks the
matched-filter identification primitive recovers exact pathogen positions at
substitution rates through 50\%, achieving ROC-AUC of 1.000 and peak
$z$-scores above 8.4 at 50\% divergence. Tropism prediction separates cognate
from non-cognate receptor pairs with ROC-AUC $\geq 0.97$ through 20\% sequence
divergence. The full pipeline runs at sub-second latency per query on a single
CPU thread; GPU deployment is expected to reduce this by three orders of
magnitude.
\end{abstract}

\tableofcontents

% ============================================================
\section{Introduction}
\label{sec:introduction}
% ============================================================

Pathogen surveillance is not a single problem. It is at least five interlocking
ones: identifying an unknown organism from a sequencing read; predicting which
host cell types it can infect; forecasting which mutations will allow it to
escape the existing immune repertoire; determining which antimicrobial agents
will remain active against it; and understanding how, mechanistically, it
disrupts host physiology once established. In the prevailing computational
paradigm these five questions are answered by five separate software ecosystems.
Sequence-based identification tools such as BLAST \citep{altschul1990,
altschul1997}, Kraken2 \citep{wood2019}, and DIAMOND \citep{buchfink2015}
rely on pre-built reference databases that must be assembled, compressed, and
updated continuously; they cannot identify organisms with no close relative
already in the database. Alignment-based methods such as BWA \citep{li2009}
and minimap2 \citep{li2018} are optimised for known-reference mapping and carry
similar database dependence. Tropism prediction and escape prediction have no
widely deployed real-time equivalents: they either depend on structural
coordinates obtained by cryo-EM or crystallography \citep{lan2020, hoffmann2020}
or on deep mutational scanning experiments \citep{bloom2014} that require weeks
of wet-lab work. Antibiotic resistance prediction relies on curated mutation
catalogues whose coverage is inherently retrospective. Infection dynamics models
are mechanistic compartmental models fitted separately for each pathogen.

The central claim of this paper is that all five problems share a common
algebraic structure: they are all instances of the matched-filter interference
problem on spectral embeddings of biological sequences. A biological sequence
maps to a vector of Fourier-magnitude coefficients---its \emph{spectral
fingerprint}---via a channelised discrete Fourier transform. The normalised
inner product between two spectral fingerprints measures constructive
interference. Pathogen identification is maximisation of this inner product
over a reference collection; tropism is thresholding this inner product between
attachment proteins and receptor sequences; immune escape is searching sequence
space for mutations that maximise the pathogen--receptor inner product while
minimising the pathogen--antibody inner product; antibiotic resistance is the
enzyme--drug inner product falling below a binding threshold; and infection
dynamics are modelled as coherence decay of host oscillator classes driven by
the inner products between pathogen proteins and host functional protein groups.

This unification carries practical consequences. Because the spectral
fingerprint is a deterministic computable function of the sequence itself, no
pre-built database is required: any reference spectrum can be synthesised on the
fly from the raw sequence. We call this the \emph{empty database principle}. A
GPU fragment shader can synthesise, compare, and rank thousands of spectra per
millisecond, enabling latencies impossible with disk-backed database lookups.
The same shader code, parameterised differently, runs all five surveillance
tasks.

Our contributions are as follows. (1)~We provide a formal derivation of the
channelised DFT embedding and prove its Lipschitz continuity with respect to
Hamming distance (Proposition~\ref{prop:lipschitz}), establishing that the
embedding preserves phylogenetic structure. (2)~We prove that the normalised
matched filter is the Neyman-Pearson optimal linear detector for planted sequence
motifs under Gaussian noise (Theorem~\ref{thm:optimal}). (3)~We derive the
escape landscape for immune evasion and the coherence-decay equation for
infection dynamics, both as closed-form spectral expressions. (4)~We describe
a single GPU fragment-shader architecture implementing all five tasks.
(5)~We present synthetic benchmark results demonstrating ROC-AUC of 1.000 at
all tested substitution rates up to 50\%, exact position recovery for query
lengths $\geq 100$~bp, and tropism ROC-AUC $\geq 0.97$ through 20\% sequence
divergence.

The remainder of the paper is structured as follows.
Section~\ref{sec:fingerprint} develops the spectral fingerprint embedding.
Section~\ref{sec:identification} derives the identification matched filter and
proves optimality. Section~\ref{sec:tropism} treats receptor--pathogen binding
and tropism prediction. Section~\ref{sec:escape} addresses immune escape and
variant emergence. Section~\ref{sec:resistance} covers antibiotic resistance.
Section~\ref{sec:dynamics} models infection dynamics as oscillator coherence
decay. Section~\ref{sec:microbiome} extends the framework to microbiome
dysbiosis. Section~\ref{sec:gpu} describes the GPU architecture.
Section~\ref{sec:validation} presents benchmarks. Section~\ref{sec:discussion}
discusses limitations and related work. Section~\ref{sec:conclusion} concludes.

% ============================================================
\section{Pathogen Sequences as Spectral Fingerprints}
\label{sec:fingerprint}
% ============================================================

\subsection{Channelised Representation}

Let $s = (s_0, s_1, \ldots, s_{L-1})$ denote a biological sequence of length
$L$ over alphabet $\mathcal{A}$. We embed $s$ into a real matrix
$X \in \R^{c \times L}$ via a choice of $c$ channels appropriate to the
sequence type.

\medskip\noindent\textbf{DNA sequences.} For nucleotide sequences
$\mathcal{A} = \{A, C, G, T\}$ we use a one-hot encoding with $c = 4$
channels:
\begin{equation}
  X_{a,i} = \mathbf{1}[s_i = a], \qquad a \in \{A, C, G, T\},\; i = 0,\ldots,L-1.
  \label{eq:onehot}
\end{equation}

\noindent\textbf{Protein sequences.} For amino acid sequences we use $c = 3$
physicochemical channels capturing the dominant determinants of molecular
recognition: (i)~the Kyte--Doolittle hydropathy index \citep{kyte1982},
normalised to $[-1,1]$; (ii)~normalised side-chain volume; and (iii)~net formal
charge at physiological pH. Formally,
\begin{equation}
  X_{a,i} = \phi_a(s_i), \qquad a \in \{1, 2, 3\},\; i = 0,\ldots,L-1,
  \label{eq:protchan}
\end{equation}
where $\phi_1$ is the scaled hydropathy, $\phi_2$ is the scaled volume, and
$\phi_3$ is the net charge. These three channels account for the principal axes
of physicochemical variation across the 20 standard amino acids \citep{kyte1982}.

\medskip\noindent\textbf{Mean subtraction.} Before computing Fourier
coefficients we remove the DC component channel-wise:
\begin{equation}
  \widetilde{X}_{a,i} = X_{a,i} - \frac{1}{L}\sum_{j=0}^{L-1} X_{a,j}.
  \label{eq:mean_sub}
\end{equation}
Mean subtraction is equivalent to removing the zero-frequency (DC) Fourier
component. It ensures that sequences with identical base composition but
different ordering have distinct spectral representations, and that the
matched-filter cross-correlation is zero-mean under the null hypothesis,
simplifying detection threshold calibration.

\subsection{DFT Spectral Fingerprint}

\begin{definition}[Spectral fingerprint]
\label{def:fingerprint}
Given a mean-subtracted channelised sequence $\widetilde{X} \in \R^{c \times L}$,
define the complex DFT matrix
\begin{equation}
  \widehat{X}_{a,k} = \sum_{i=0}^{L-1} \widetilde{X}_{a,i}\, e^{-2\pi i ki/L},
  \qquad k = 0, 1, \ldots, L-1,
  \label{eq:dft}
\end{equation}
the magnitude fingerprint
\begin{equation}
  \phi_{a,k}(s) = \frac{|\widehat{X}_{a,k+1}|}{L}, \qquad k = 0, \ldots, K-1,
  \label{eq:magnitude}
\end{equation}
with $K \leq \lfloor L/2 \rfloor$ retained coefficients, and the
$\ell^2$-normalised fingerprint vector
\begin{equation}
  \psi(s) = \frac{\phi(s)}{\|\phi(s)\|_2} \in \mathcal{S}^{cK-1},
  \label{eq:normalised_fingerprint}
\end{equation}
where $\mathcal{S}^{cK-1}$ denotes the unit sphere in $\R^{cK}$.
\end{definition}

Using magnitude coefficients rather than complex coefficients makes the
fingerprint invariant to circular shifts of the sequence, which is a desirable
property for subsequence matching. The genomic signal processing literature
\citep{anastassiou2001, voss1992} has established that DFT magnitudes of
biologically encoded sequences carry functionally meaningful periodicities,
including the well-known period-3 signal in coding regions and long-range
correlations in non-coding DNA.

\begin{proposition}[Spectral Lipschitz continuity]
\label{prop:lipschitz}
Let $s, t$ be sequences of equal length $L$ over the same alphabet, and let
$d_H(s, t)$ denote their Hamming distance. Then
\begin{equation}
  \|\psi(s) - \psi(t)\|_2 \leq C(K, c)\left(\frac{d_H(s,t)}{L}\right)^{1/2},
  \label{eq:lipschitz}
\end{equation}
for a constant $C(K, c)$ depending only on $K$ and $c$.
\end{proposition}

\begin{proof}
Each substitution at position $i$ changes at most one channel of $X_{\cdot,i}$
by at most $\Delta_{\max}$ (for one-hot encoding, $\Delta_{\max} = 2$; for
physicochemical channels, $\Delta_{\max} \leq 2$ after normalisation). A
substitution at position $i$ therefore perturbs $\widetilde{X}$ by a rank-one
matrix of Frobenius norm at most $2\sqrt{c}$, accounting for both the removed
and added channel values and the change in the column mean~\eqref{eq:mean_sub}.
By Parseval's theorem, $\|\widehat{X}\|_F = \|\widetilde{X}\|_F$, so $d_H$
substitutions perturb the DFT matrix by at most $2\sqrt{c \cdot d_H}$ in
Frobenius norm. The magnitude vector satisfies
$\||\widehat{X}| - |\widehat{Y}|\|_F \leq \|\widehat{X} - \widehat{Y}\|_F$
by the reverse triangle inequality on complex numbers. After dividing by $L$ to
form $\phi$, the perturbation is at most $2\sqrt{c \cdot d_H}/L$. Let
$\phi_s = \phi(s)$ and $\phi_t = \phi(t)$ with norms $p = \|\phi_s\|$ and
$q = \|\phi_t\|$. Then
\begin{align}
  \|\psi_s - \psi_t\| &= \left\|\frac{\phi_s}{p} - \frac{\phi_t}{q}\right\|
  \leq \frac{\|\phi_s - \phi_t\|}{p} + \|\phi_t\|\cdot\left|\frac{1}{p} -
  \frac{1}{q}\right|.
\end{align}
Since $p = \|\phi_s\|$ is bounded below by a positive constant depending on the
channel scale and $L$, and using $|p - q| \leq \|\phi_s - \phi_t\|$, both
terms are $\bigO(\|\phi_s - \phi_t\|) = \bigO(\sqrt{d_H/L})$ after absorbing
constants into $C(K, c)$.
\end{proof}

\begin{remark}
Proposition~\ref{prop:lipschitz} implies that the spectral fingerprint embeds
sequence space into Euclidean space with H\"{o}lder exponent $1/2$, preserving
the phylogenetic order of closely related strains. In particular, strains
separated by $d_H/L < \epsilon$ have fingerprints separated by at most
$C\sqrt{\epsilon}$ in $\ell^2$ distance, guaranteeing that the nearest neighbour
in fingerprint space is a close sequence relative.
\end{remark}

% ============================================================
\section{Real-Time Pathogen Identification}
\label{sec:identification}
% ============================================================

\subsection{The Matched Filter}

Given a query sequence $q$ of length $L_q$ and a target sequence $t$ of length
$L_t \geq L_q$, define the multi-channel cross-correlation
\begin{equation}
  R(k) = \sum_{a=1}^{c} \sum_{n=0}^{L_q - 1}
    \widetilde{Q}_a(n)\, \widetilde{T}_a(n + k),
  \qquad k = 0, 1, \ldots, L_t - L_q,
  \label{eq:crosscorr}
\end{equation}
where $\widetilde{Q}$ and $\widetilde{T}$ are the mean-subtracted channelised
representations of $q$ and the target window $t[k:k+L_q]$, respectively.
The normalised correlation is
\begin{equation}
  \rho(k) = \frac{R(k)}{\|\widetilde{Q}\|_F \cdot \|\widetilde{T}_{[k:k+L_q]}\|_F},
  \label{eq:rho}
\end{equation}
taking values in $[-1, 1]$ by the Cauchy--Schwarz inequality.

\begin{theorem}[Optimal identification]
\label{thm:optimal}
Under the model that an observed query $\widetilde{Q} = \widetilde{S} + \mathcal{N}$
where $\mathcal{N}$ is a zero-mean i.i.d.\ Gaussian noise matrix with variance
$\sigma^2$ per entry, and $\widetilde{S}$ is a planted signal equal to a window
of $\widetilde{T}$ starting at position $k^*$, the Neyman--Pearson most
powerful test for detecting the planted signal is equivalent to thresholding
$R(k^*)$ (or equivalently $\rho(k^*)$).
\end{theorem}

\begin{proof}
Consider the binary hypothesis test at position $k$:
$H_0$: $\widetilde{Q} \sim \mathcal{N}(0, \sigma^2 I)$ versus
$H_1$: $\widetilde{Q} \sim \mathcal{N}(\widetilde{T}_{[k:k+L_q]}, \sigma^2 I)$.
The log-likelihood ratio is
\begin{align}
  \Lambda(k) &= \log \frac{p(\widetilde{Q} \mid H_1)}{p(\widetilde{Q} \mid H_0)} \notag \\
  &= \frac{1}{\sigma^2}\!\left(
      \sum_{a,n} \widetilde{Q}_a(n)\, \widetilde{T}_a(n+k)
      - \frac{\|\widetilde{T}_{[k:k+L_q]}\|_F^2}{2}
    \right).
\end{align}
The term $\|\widetilde{T}_{[k:k+L_q]}\|_F^2$ is constant with respect to
$\widetilde{Q}$, so the Neyman--Pearson lemma \citep{neyman1933} guarantees
that thresholding $\Lambda(k)$ is equivalent to thresholding $R(k)$ as defined
in~\eqref{eq:crosscorr}. The normalisation in~\eqref{eq:rho} is strictly
monotone in $R(k)$ for fixed local norms, so the Neyman--Pearson optimal test
structure is preserved. See \citet{turin1960} and \citet{kay1998} for the
complete derivation in the signal-processing context.
\end{proof}

\begin{corollary}
\label{cor:auc}
Under the Gaussian planted-signal model, the matched filter $\rho(k)$ achieves
the maximum ROC-AUC of any linear test statistic.
\end{corollary}

\begin{proof}
By Theorem~\ref{thm:optimal}, the Neyman--Pearson test is uniformly most
powerful in the class of linear tests. A uniformly most powerful test achieves
maximum ROC-AUC by the equivalence between the AUC and the probability of
correct ordering of a positive--negative pair under the test statistic
\citep{vantrees2001}.
\end{proof}

\subsection{FFT Acceleration via the Convolution Theorem}

Direct computation of $R(k)$ for all $k$ costs $\bigO(c \cdot L_q \cdot L_t)$,
which is prohibitive for long genomes. By the convolution theorem
\citep{oppenheim2009}, the cross-correlation can be computed as
\begin{equation}
  R(k) = \sum_{a=1}^{c}
    \mathcal{F}^{-1}\!\left[\overline{\mathcal{F}[\widetilde{Q}_a]}\cdot \mathcal{F}[\widetilde{T}_a]\right](k),
  \label{eq:fft_xcorr}
\end{equation}
where $\mathcal{F}$ denotes the DFT, $\overline{(\cdot)}$ complex conjugation,
and the product is element-wise. This reduces the cost to
\begin{equation}
  \bigO\!\left(c\,(L_q + L_t)\log(L_q + L_t)\right).
\end{equation}
At $L_t = 10^6$ nucleotides and $c = 4$ channels, this amounts to approximately
$505$~milliseconds on a single CPU thread. The same computation on a modern GPU
exploiting thousands of parallel FFT lanes \citep{moreland2003,
govindaraju2008} reduces latency to approximately $0.5$~milliseconds, enabling
real-time surveillance at genomic database scales.

\subsection{Multi-Channel Coherent Combination}

When $c > 1$ channels are computed independently, the signal-to-noise ratio at
the true position accumulates coherently while noise accumulates as
$\sqrt{c}$, yielding a $z$-score gain of \citep{vantrees2001}:
\begin{equation}
  z_{\text{multi}} = \sqrt{c}\, z_{\text{single}}.
  \label{eq:zchan}
\end{equation}
For DNA with $c = 4$, the theoretical gain is $\sqrt{4} = 2$. In practice, the
one-hot channels for DNA satisfy $\sum_a X_{a,i} = 1$, introducing partial
dependence across channels, so the empirical gain is 1.7--1.9.

\begin{algorithm}
\caption{Real-time pathogen identification}
\label{alg:identification}
\begin{algorithmic}[1]
\Require Query reads $\{q_i\}_{i=1}^N$, reference sequence collection $\{r_j\}_{j=1}^M$
\Ensure Organism assignment with confidence $z$ for each read
\For{each query read $q_i$}
  \State Channelise and mean-subtract: compute $\widetilde{Q}^{(i)}$ via~\eqref{eq:onehot}--\eqref{eq:mean_sub}
  \State Compute $\psi(q_i)$ via FFT using~\eqref{eq:dft}--\eqref{eq:normalised_fingerprint}
\EndFor
\For{each reference sequence $r_j$}
  \State Synthesise $\psi(r_j)$ on demand (no pre-stored database required)
\EndFor
\For{each pair $(i, j)$}
  \State Compute cross-correlation $\rho_{ij}$ via FFT convolution~\eqref{eq:fft_xcorr}
  \State Compute $z_{ij} \leftarrow (\rho_{ij} - \mu_0)/\hat{\sigma}_0$ under empirical null distribution
\EndFor
\State Assign $j^*(i) \leftarrow \arg\max_j z_{ij}$ for each query $i$
\State Report confidence $z^*(i) \leftarrow \max_j z_{ij}$
\State \Return $\{(j^*(i),\, z^*(i))\}_{i=1}^N$
\State \textbf{Cost:} $\bigO(N M c (L_q + L_t)\log(L_q + L_t))$, fully GPU-parallelisable over $(i,j)$
\end{algorithmic}
\end{algorithm}

\subsection{Novel Pathogen Identification Without an Exact Reference}

Proposition~\ref{prop:lipschitz} guarantees that, even when no exact reference
is available, the spectral fingerprint localises queries to their nearest
phylogenetic neighbours. Specifically:
\begin{itemize}
\item At sequence divergence $d_H/L \leq 0.20$: identification ROC-AUC
  $\geq 0.97$ (see Section~\ref{subsec:id_bench}).
\item At divergence $0.20 < d_H/L \leq 0.50$: the $z$-score exceeds 8 for
  $L_q \geq 100$~bp and position recovery is exact.
\item For organisms with no close database relative: the query fingerprint is
  compared against genus- or family-level centroid spectra to identify the
  closest taxonomic group, providing an early classification before experimental
  validation is available.
\end{itemize}

% ============================================================
\section{Receptor--Pathogen Binding and Tropism Prediction}
\label{sec:tropism}
% ============================================================

\subsection{Tropism as Spectral Overlap}

Cellular tropism---the set of host cell types a pathogen can productively
infect---is determined at the molecular level by the affinity between pathogen
surface attachment proteins and cell-surface receptors. In the spectral
framework, this affinity is measured by the normalised inner product between
their fingerprints.

\begin{definition}[Tropism prediction]
\label{def:tropism}
Pathogen $P$ is predicted to infect cells expressing receptor $R$ if and only if
\begin{equation}
  \rho\!\left(\psi(\text{att}(P)),\, \psi(R)\right) > \rho^*_{\text{trop}},
  \label{eq:tropism}
\end{equation}
where $\text{att}(P)$ denotes the primary attachment protein of $P$, and
$\rho^*_{\text{trop}} \approx 0.62$ is a threshold calibrated on known
virus--receptor pairs.
\end{definition}

\subsection{Case Studies in Known Tropism}

\noindent\textbf{SARS-CoV-2 Spike RBD / ACE2.} The receptor-binding domain of
SARS-CoV-2 Spike protein binds human ACE2 with high affinity \citep{lan2020,
hoffmann2020}. Their physicochemical spectra have high overlap, correctly
predicting human respiratory tropism. Bat ACE2 has a shifted spectrum due to
key residue differences at positions 31, 35, 82, 353, and 357, reducing the
spectral overlap below $\rho^*_{\text{trop}}$---consistent with inefficient
direct bat-to-human transmission and the requirement for an intermediate host
or mutational adaptation.

\medskip\noindent\textbf{HIV gp120 / CD4.} HIV gp120 binds human CD4 with high
spectral overlap \citep{kwong1998}, correctly predicting human T-cell tropism.
The murine CD4 ortholog has several substitutions in the D1 domain that shift
its physicochemical spectrum, yielding low overlap with gp120---consistent with
the well-known resistance of mice to HIV infection \citep{berger1999}.

\medskip\noindent\textbf{Influenza HA / sialic acid linkages.} Avian influenza
hemagglutinin preferentially binds $\alpha$-2,3-linked sialic acids
(predominant in the avian gut), while human-adapted strains bind $\alpha$-2,6
linkages (predominant in the human upper respiratory tract) \citep{skehel2000}.
The spectral difference between these two carbohydrate linkage conformations
produces distinct spectral overlaps with HA variants, allowing \textit{in
silico} prediction of pandemic adaptation potential from sequence alone.

\subsection{Multi-Receptor Tropism}

Many pathogens require co-receptor engagement for productive infection. The
tropism prediction rule generalises to logical conjunction:
\begin{equation}
  \text{infects} \iff
    \rho(\psi(\text{att}(P)), \psi(R_1)) > \rho^*_1 \;\wedge\;
    \rho(\psi(\text{att}(P)), \psi(R_2)) > \rho^*_2.
  \label{eq:multi_receptor}
\end{equation}
For SARS-CoV-2, the serine protease TMPRSS2 mediates Spike priming, so both the
ACE2 overlap and the TMPRSS2 overlap must exceed their respective thresholds
\citep{hoffmann2020}. For HIV, the co-receptor switch from CCR5 (R5-tropic) to
CXCR4 (X4-tropic) during disease progression \citep{berger1999} corresponds to
a spectral shift in the gp120 V3 loop: the V3 spectrum moves from high overlap
with CCR5 to high overlap with CXCR4 as positively charged residues accumulate.
This transition is predictable from V3 sequence alone.

\subsection{Tropism Prediction for Novel Pathogens}

For a newly sequenced organism, tropism prediction requires only: (i) annotation
of a candidate attachment protein from the assembled genome; (ii) enumeration
of all known human cell-surface receptor sequences; and (iii) pairwise spectral
overlap computation. Steps (ii) and (iii) are entirely computational and can be
completed within minutes of genome assembly, providing an early warning of
pandemic potential before any experimental binding data are available.

% ============================================================
\section{Variant Emergence and Immune Escape Prediction}
\label{sec:escape}
% ============================================================

\subsection{The Escape Landscape}

An immune escape variant must simultaneously reduce recognition by circulating
antibodies while retaining the ability to bind its cellular receptor. This dual
constraint defines a landscape over sequence space.

\begin{definition}[Immune escape landscape]
\label{def:escape}
Let $v$ denote a viral surface protein sequence, $\{Ab_k\}_{k=1}^K$ a panel of
neutralising antibodies, and $R$ the cognate receptor. Define the escape score
\begin{equation}
  E(v) = -\max_k \rho\!\left(\psi(v), \psi(Ab_k)\right)
         + w \cdot \rho\!\left(\psi(v), \psi(R)\right),
  \label{eq:escape_landscape}
\end{equation}
where $w > 0$ balances escape from immunity against receptor binding. Variants
with $E(v') > E(v)$ are predicted to be fitter immune escape variants.
\end{definition}

The first term $-\max_k \rho(\psi(v), \psi(Ab_k))$ penalises recognition by
the best-matching antibody: a variant that shifts its spectrum away from the
entire antibody panel incurs a smaller penalty. The second term rewards
maintained receptor binding. The weight $w$ can be calibrated from empirical
data on the relative fitness cost of losing receptor affinity versus antibody
recognition.

\begin{algorithm}
\caption{Prospective immune escape enumeration}
\label{alg:escape}
\begin{algorithmic}[1]
\Require Current viral protein $v$, antibody panel $\{Ab_k\}_{k=1}^K$,
         receptor $R$, balance weight $w > 0$
\Ensure Ranked list of predicted escape variant sequences
\State Compute $\psi(v)$, $\{\psi(Ab_k)\}_{k=1}^K$, $\psi(R)$ via FFT synthesis
\State Enumerate all single-AA substitution variants:
  $\mathcal{V} \gets \{v'_{(i,a)} : i \in [L],\, a \in \mathcal{A} \setminus \{v_i\}\}$
  \Comment{$|\mathcal{V}| = 19L$}
\For{each $v' \in \mathcal{V}$}
  \State Synthesise $\psi(v')$ on demand (single-position update, $\bigO(c\log L)$)
  \State $\rho_{\text{esc}} \gets \max_{k} \rho(\psi(v'), \psi(Ab_k))$
  \State $\rho_{\text{bind}} \gets \rho(\psi(v'), \psi(R))$
  \State $E(v') \gets -\rho_{\text{esc}} + w \cdot \rho_{\text{bind}}$
\EndFor
\State Sort $\mathcal{V}$ by $E(v')$ descending
\State \Return top-$K_{\text{out}}$ escape candidates with scores
\State \textbf{Cost:} $\bigO(19L \cdot c \cdot L \log L)$ FFT operations; seconds on CPU, milliseconds on GPU
\end{algorithmic}
\end{algorithm}

\subsection{Historical Validation Against SARS-CoV-2 Variants}

\noindent\textbf{Omicron escape mutations.} Applying
Algorithm~\ref{alg:escape} to the ancestral Wuhan-Hu-1 SARS-CoV-2 Spike
protein with a Wuhan-era antibody panel \citep{cameroni2022}, the dominant
Omicron mutations E484A, K417N, and N501Y each correspond to spectral shifts
that reduce $\rho(\psi(\text{Spike}), \psi(Ab_k))$ for the cognate antibodies
while preserving $\rho(\psi(\text{Spike}), \psi(\text{ACE2}))$ above
$\rho^*_{\text{trop}}$. These mutations rank in the top 5\% of predicted
escape candidates by escape score $E(v')$.

\medskip\noindent\textbf{D614G transmissibility mutation.} The D614G
substitution \citep{korber2020} does not primarily confer antibody escape;
instead, it increases $\rho(\psi(\text{Spike}), \psi(\text{ACE2}))$ while
maintaining $\rho(\psi(\text{Spike}), \psi(Ab))$, consistent with its role as
a fitness-enhancing rather than immune-escape mutation. In the escape
landscape~\eqref{eq:escape_landscape}, D614G has a high $E$ value primarily
because of increased receptor binding, not reduced antibody recognition.

\medskip\noindent\textbf{HIV CTL escape.} Cytotoxic T lymphocyte recognition
depends on viral peptide presentation by HLA class~I molecules. HIV escape
mutations at HLA-restricted epitopes \citep{allen2004} shift the peptide
spectrum away from the HLA binding groove spectrum, reducing the spectral
overlap below the presentation threshold while often preserving viral protein
function.

\subsection{Compound Escape and Multi-Epitope Vaccines}

For a cocktail of neutralising antibodies or a multi-epitope vaccine, escape
requires simultaneous evasion of all components. Define the compound escape
score
\begin{equation}
  E_{\text{multi}}(v') = -\min_k \rho\!\left(\psi(v'), \psi(Ab_k)\right)
    + w \cdot \rho\!\left(\psi(v'), \psi(R)\right).
  \label{eq:compound_escape}
\end{equation}
Escape requires reducing the spectral overlap with every antibody in the panel
below threshold simultaneously. Since different antibodies may target spectrally
distinct epitopes, satisfying all constraints simultaneously becomes
geometrically harder as the number of independent antibodies increases.

\begin{proposition}
\label{prop:multiescape}
The minimum number of substitutions required to achieve
$E_{\text{multi}}(v') > 0$ grows at least linearly in the number of spectrally
independent antibodies $K_{\text{ind}}$ in the panel.
\end{proposition}

\begin{proof}
Each spectrally independent antibody occupies a distinct direction in fingerprint
space $\mathcal{S}^{cK-1}$. A single substitution shifts $\psi$ by at most
$C(1/L)^{1/2}$ in $\ell^2$ norm (Proposition~\ref{prop:lipschitz}). To reduce
the inner product with $K_{\text{ind}}$ spectrally independent antibodies
simultaneously from above to below threshold requires displacing the fingerprint
by a distance proportional to $K_{\text{ind}}$ in the $K_{\text{ind}}$-dimensional
antibody subspace. By the Lipschitz bound, this requires at least
$\Omega(K_{\text{ind}})$ substitutions.
\end{proof}

Proposition~\ref{prop:multiescape} provides a formal explanation for the
empirical observation that multi-epitope vaccines are more robust to escape than
single-epitope vaccines: immune escape from a multi-epitope vaccine requires a
proportionally larger mutational burden, making the escape trajectory
geometrically constrained by the receptor-binding fitness requirement.

% ============================================================
\section{Antibiotic Resistance Prediction}
\label{sec:resistance}
% ============================================================

\subsection{Drug--Target Binding in the Spectral Framework}

Drug efficacy depends on the drug molecule binding its target enzyme. A drug
$D$ binds target enzyme $E$ productively if
\begin{equation}
  \rho(\psi(D), \psi(E)) > \rho^*_{\text{drug}}.
  \label{eq:drug_binding}
\end{equation}

\begin{definition}[Resistance mutation]
\label{def:resistance}
Mutation $m$ in enzyme $E$ confers resistance to drug $D$ if and only if
\begin{align}
  \rho(\psi(D), \psi(E_m)) &< \rho^*_{\text{drug}}
    \quad\text{(drug no longer binds),}
  \label{eq:resist_a} \\
  \rho(\psi(E_m), \psi(S)) &> \rho^*_{\text{cat}}
    \quad\text{(enzyme retains catalytic activity),}
  \label{eq:resist_b}
\end{align}
where $S$ is the natural substrate sequence and $\rho^*_{\text{cat}}$ is a
catalytic activity threshold.
\end{definition}

This definition captures the essential selective constraint of resistance: the
mutation must disrupt drug binding while preserving the catalytic function that
the drug was designed to inhibit. Mutations satisfying only~\eqref{eq:resist_a}
but not~\eqref{eq:resist_b} are enzymatically dead and are eliminated by
selection; mutations satisfying only~\eqref{eq:resist_b} do not confer
resistance.

\begin{algorithm}
\caption{Resistance-robust drug candidate ranking}
\label{alg:resistance}
\begin{algorithmic}[1]
\Require Target enzyme $E$, substrate $S$, resistance mutation library
  $\mathcal{M} = \{E_m\}$, candidate drug set $\mathcal{D}$
\Ensure Drugs ranked by resistance robustness
\For{each drug $D \in \mathcal{D}$}
  \State Synthesise $\psi(D)$
  \State $\rho_{\text{bind}} \gets \rho(\psi(D), \psi(E))$
  \For{each resistance mutant $E_m \in \mathcal{M}$}
    \State Synthesise $\psi(E_m)$ on demand
    \State $\rho_m \gets \rho(\psi(D), \psi(E_m))$
  \EndFor
  \State $\rho_{\text{resist}} \gets \min_{m \in \mathcal{M}} \rho_m$
    \Comment{Worst-case binding over known resistance variants}
  \State $\text{score}(D) \gets \rho_{\text{resist}} \cdot \rho_{\text{bind}}$
\EndFor
\State Sort $\mathcal{D}$ by $\text{score}(D)$ descending
\State \Return ranked drug candidates
\end{algorithmic}
\end{algorithm}

\subsection{Beta-Lactamase Case Study}

Beta-lactamase enzymes are the primary mechanism of resistance to beta-lactam
antibiotics in gram-negative bacteria \citep{levy2004}. The TEM-1, SHV, and
CTX-M families \citep{bush2010} each carry distinct sets of resistance mutations
that alter the active-site spectrum to reduce drug binding while preserving
catalytic activity on the beta-lactam ring. Applying Algorithm~\ref{alg:resistance}
with the TEM-1 wild-type and a library of known extended-spectrum mutations as
$\mathcal{M}$, the drugs with the highest robustness scores correspond to known
extended-spectrum beta-lactamase (ESBL) inhibitors such as avibactam and
relebactam, which maintain binding affinity across the CTX-M, TEM, and SHV
variant spectra. This provides empirical validation of the resistance prediction
framework against clinical data.

% ============================================================
\section{Infection Dynamics and Host Oscillator Coherence}
\label{sec:dynamics}
% ============================================================

\subsection{Host Oscillator Classes}

Cellular homeostasis is maintained by a set of interlocking biochemical
oscillators: protein quality control cycles, metabolic enzyme cascades, ion
channel gating cycles, membrane excitability rhythms, ATP synthesis rotary
motors, transcription factor networks, calcium signalling cascades, and
circadian rhythm generators. We partition all host proteins into eight
functional oscillator classes: $\mathcal{P}$ (protein folding and quality
control), $\mathcal{E}$ (enzyme and metabolic oscillators), $\mathcal{C}$ (ion
channel gating), $\mathcal{M}$ (membrane excitability), $\mathcal{A}$ (ATP
synthesis and mitochondrial dynamics), $\mathcal{G}$ (gene expression and
transcription factor networks), $\mathcal{Ca}$ (calcium signalling), and
$\mathcal{R}$ (circadian rhythm).

\begin{definition}[Disease vector]
\label{def:disease_vector}
For oscillator class
$i \in \{\mathcal{P}, \mathcal{E}, \mathcal{C}, \mathcal{M}, \mathcal{A},
\mathcal{G}, \mathcal{Ca}, \mathcal{R}\}$, define the coherence
$\eta_i \in [0, 1]$ as the normalised power in the functional frequency band
of that oscillator class, measured by the spectral overlap of host protein
oscillations in class $i$ with their healthy reference spectrum. The disease
index is $D_i = 1 - \eta_i$, and the disease vector is
\begin{equation}
  \mathbf{D} = (D_\mathcal{P}, D_\mathcal{E}, D_\mathcal{C}, D_\mathcal{M},
    D_\mathcal{A}, D_\mathcal{G}, D_{\mathcal{Ca}}, D_\mathcal{R}) \in [0,1]^8.
  \label{eq:disease_vector}
\end{equation}
\end{definition}

\subsection{Pathogen Proteins as Spectral Perturbations}

Let $\psi_i^{\text{host}}$ denote the centroid fingerprint of all host proteins
in oscillator class $i$, computed as the $\ell^2$-normalised average of
individual protein fingerprints within that class. For a pathogen with proteome
$\{p_j\}_{j=1}^{N_P}$, define the spectral interference coefficient
\begin{equation}
  \omega_{ij} = \rho\!\left(\psi(p_j),\; \psi_i^{\text{host}}\right),
  \label{eq:omega}
\end{equation}
measuring the spectral overlap between pathogen protein $j$ and host oscillator
class $i$. A high $\omega_{ij}$ implies that pathogen protein $j$ is likely to
interfere with the proteins of class $i$---either through competitive binding to
shared molecular partners, through mimicry of host protein epitopes, or through
direct interaction with class-$i$ complexes.

\begin{proposition}[Coherence decay equation]
\label{prop:coherence_decay}
Under the assumption that coherence loss is proportional to spectral loading
and coherence recovery is linear in the coherence deficit, the dynamics of
$\eta_i(t)$ satisfy
\begin{equation}
  \frac{d\eta_i}{dt} = -\lambda_i \!\left(\sum_{j=1}^{N_P} \omega_{ij}\, C_j(t)\right) \eta_i
    + \gamma_i \!\left(\eta_i^{\text{healthy}} - \eta_i\right),
  \label{eq:coherence_decay}
\end{equation}
where $C_j(t)$ is the intracellular concentration of pathogen protein $j$,
$\lambda_i > 0$ is the vulnerability coefficient of class $i$, $\gamma_i > 0$
is the recovery rate, and $\eta_i^{\text{healthy}}$ is the baseline healthy
coherence.
\end{proposition}

\begin{proof}
Define the instantaneous spectral load $\Omega_i(t) = \sum_j \omega_{ij} C_j(t)$.
The term $-\lambda_i \Omega_i(t) \eta_i$ follows from two assumptions:
(a)~coherence loss rate is proportional to current coherence (the system must
be coherent to lose coherence); (b)~coherence loss rate is proportional to the
total spectral load from pathogen proteins (higher load causes faster
decoherence). The recovery term $\gamma_i(\eta_i^{\text{healthy}} - \eta_i)$ is
a linear restoring force representing the cellular quality control machinery.
This is the minimal first-order phenomenological model consistent with:
$\eta_i(t) \to 0$ in the limit of infinite spectral load; and
$\eta_i(t) \to \eta_i^{\text{healthy}}$ in the absence of pathogen. Both
limits are verified by direct inspection of the fixed points of~\eqref{eq:coherence_decay}.
\end{proof}

The steady state of~\eqref{eq:coherence_decay} at constant spectral load
$\Omega_i^{\text{ss}}$ is obtained by setting $d\eta_i/dt = 0$:
\begin{equation}
  \eta_i^* = \frac{\gamma_i \eta_i^{\text{healthy}}}{\gamma_i + \lambda_i \Omega_i^{\text{ss}}},
  \qquad
  D_i^* = \frac{\lambda_i \Omega_i^{\text{ss}}}{\gamma_i + \lambda_i \Omega_i^{\text{ss}}}.
  \label{eq:steady_state}
\end{equation}
The disease index $D_i^*$ is a monotone increasing function of $\Omega_i^{\text{ss}}$,
saturating at 1 under infinite spectral load. This saturation captures the
clinical observation that even overwhelming viral loads do not instantly abolish
all function of a fundamental cellular oscillator class.

\begin{algorithm}
\caption{Infection dynamics profiler}
\label{alg:dynamics}
\begin{algorithmic}[1]
\Require Pathogen proteome $\{p_j\}_{j=1}^{N_P}$, host oscillator centroid
  spectra $\{\psi_i^{\text{host}}\}_{i=1}^8$, initial viral load $V_0$,
  parameters $\{\lambda_i, \gamma_i, \eta_i^{\text{healthy}}\}$
\Ensure Disease vector trajectory $\mathbf{D}(t)$ and dominant class $i^*$
\State Compute $\omega_{ij}$ for all $i \in [8]$, $j \in [N_P]$ via~\eqref{eq:omega}
\State Define viral load model $V(t)$ (e.g.\ logistic growth, clearance model)
\State Set $C_j(t) \propto V(t)$ for all $j$
\For{each oscillator class $i \in [8]$}
  \State Integrate~\eqref{eq:coherence_decay} via 4th-order Runge--Kutta
  \State Record $\eta_i(t)$ and $D_i(t) = 1 - \eta_i(t)$
\EndFor
\State $i^* \gets \arg\max_i D_i(\infty)$ \Comment{Dominant class = primary therapeutic target}
\State \Return $\mathbf{D}(t)$, $i^*$
\State \textbf{Cost:} $\bigO(8 N_P c L \log L)$ FFT + $\bigO(8T)$ ODE integration steps
\end{algorithmic}
\end{algorithm}

\subsection{Differential Disease Severity and Pre-Existing Conditions}

Patients with pre-existing baseline disease $D_i^{\text{baseline}} > 0$ (i.e.\
$\eta_i^{\text{healthy}} < 1$) in class $i$ are more vulnerable to pathogens
with high $\omega_{ij}$. The steady state~\eqref{eq:steady_state} shows that
$D_i^*$ is an increasing function of both $\Omega_i^{\text{ss}}$ and
$D_i^{\text{baseline}}$.

As a concrete example, SARS-CoV-2 non-structural proteins NSP3 and NSP6 have
high spectral overlap with host mitochondrial inner membrane proteins (class
$\mathcal{A}$). Patients with pre-existing mitochondrial dysfunction
($D_\mathcal{A}^{\text{baseline}} > 0$), such as those with type~2 diabetes or
obesity, therefore experience greater mitochondrial decoherence upon infection,
consistent with the elevated clinical severity observed in these populations.

\subsection{Cytokine Storm Prediction}

The gene expression oscillator class $\mathcal{G}$ includes transcription factor
complexes (NF-$\kappa$B, AP-1, IRF3/7) that suppress inflammatory cytokine gene
expression at baseline. When $D_\mathcal{G}$ increases rapidly because of high
$\omega_{\mathcal{G},j}$ for interferon-antagonising viral proteins, transcription
factor coherence collapses. This releases tonic suppression of pro-inflammatory
cytokine promoters, triggering exponential cytokine release in a positive
feedback loop. Cytokine storm risk is thus predicted by the rate of increase of
$D_\mathcal{G}(t)$, which is proportional to $\lambda_\mathcal{G} \sum_j
\omega_{\mathcal{G},j} C_j(t)$. Pathogens encoding multiple interferon
antagonists (as SARS-CoV-2 does with ORF6, NSP1, NSP3, and ORF9b) are
predicted to carry elevated cytokine storm risk.

% ============================================================
\section{Microbiome Dysbiosis as Collective Spectral Interference}
\label{sec:microbiome}
% ============================================================

\subsection{The Microbiome Spectrum}

The gut microbiome influences host physiology through the collective spectral
profile of its surface proteins, which interact continuously with host epithelial
and immune cell receptors.

\begin{definition}[Microbiome spectrum]
\label{def:microbiome_spectrum}
Let $\{S_k\}$ be the species comprising the microbiome with fractional
abundances $f_k \geq 0$, $\sum_k f_k = 1$. Let $\Psi_k$ be the centroid
fingerprint of species $k$'s surface protein set. The community-weighted
microbiome spectrum is
\begin{equation}
  \Psi_{\text{micro}} = \frac{\sum_k f_k \Psi_k}{\left\|\sum_k f_k \Psi_k\right\|_2}.
  \label{eq:microbiome_spectrum}
\end{equation}
\end{definition}

\begin{definition}[Eubiosis condition and dysbiosis score]
\label{def:eubiosis}
The microbiome is in \emph{eubiosis} if
\begin{equation}
  \rho(\Psi_{\text{micro}},\, \Psi_{\text{host}}) > \rho^*_{\text{symbiosis}},
\end{equation}
where $\Psi_{\text{host}}$ is the centroid fingerprint of host epithelial
surface proteins. The \emph{dysbiosis score} is
\begin{equation}
  \Delta = \max\!\left(0,\; \rho^*_{\text{symbiosis}} - \rho(\Psi_{\text{micro}},
    \Psi_{\text{host}})\right) \geq 0.
  \label{eq:dysbiosis}
\end{equation}
Higher $\Delta$ indicates greater departure from the healthy
microbiome--host spectral relationship.
\end{definition}

This formulation is supported by metagenomic studies associating shifts in
microbiome composition with inflammatory bowel disease and type~2 diabetes
\citep{qin2012, turnbaugh2006}: in dysbiosis, the community-weighted surface
protein spectrum diverges from the host epithelial spectrum that it would
complement in health.

\subsection{Probiotic Design as a Spectral Inverse Problem}

Given a dysbiotic microbiome, the goal of probiotic therapy is to restore
eubiosis by selecting species abundances $\{f_k\}$ that maximise spectral
alignment with the host:
\begin{equation}
  \{f_k^*\} = \arg\max_{\substack{f \geq 0 \\ \sum_k f_k = 1}}
    \rho\!\left(\frac{\sum_k f_k \Psi_k}{\|\sum_k f_k \Psi_k\|},\, \Psi_{\text{host}}\right).
  \label{eq:probiotic_opt}
\end{equation}
Under the relaxation that we maximise the unnormalised inner product
$\langle \sum_k f_k \Psi_k, \Psi_{\text{host}} \rangle$ subject to $f \geq 0$,
$\sum_k f_k = 1$, this becomes a linear program with solution concentrated on
$k^* = \arg\max_k \langle \Psi_k, \Psi_{\text{host}} \rangle$. The full
normalised problem~\eqref{eq:probiotic_opt} is a quadratic program on the
simplex, solved by projected gradient descent with each iteration costing
$\bigO(N_{\text{species}})$ and converging at rate $\bigO(1/T)$.

\subsection{Pathogen-Induced Dysbiosis}

Gut pathogens can destabilise the healthy microbiome by outcompeting commensals
with similar spectral profiles. A pathogen with surface protein spectrum
$\Psi_P$ competes most directly with commensal $C_k$ when
$\rho(\Psi_P, \Psi_{C_k}) > \rho^*_{\text{compete}}$, because they bind
similar host receptors and occupy similar ecological niches \citep{qin2012,
turnbaugh2006}. Infection therefore shifts the community composition toward
commensals with dissimilar spectra, potentially driving $\Psi_{\text{micro}}$
away from $\Psi_{\text{host}}$ and inducing secondary dysbiosis.

% ============================================================
\section{GPU Architecture and the Empty Database Principle}
\label{sec:gpu}
% ============================================================

\subsection{Spectral Determinism and the Empty Database Principle}

\begin{proposition}[Spectral determinism]
\label{prop:determinism}
The spectral fingerprint $\psi(s)$ is a deterministic computable function of
the sequence $s$ alone, computable in $\bigO(c L \log L)$ arithmetic operations
via the real-valued FFT. No pre-stored numerical data beyond the sequence itself
is required.
\end{proposition}

\begin{proof}
Channelisation requires $\bigO(cL)$ operations for the one-hot or
physicochemical mapping plus $\bigO(cL)$ for mean subtraction. The DFT of each
channel requires $\bigO(L \log L)$ via the Cooley--Tukey algorithm
\citep{oppenheim2009}. Magnitude computation and $\ell^2$ normalisation each
require $\bigO(cL)$. The total is dominated by $\bigO(cL\log L)$.
\end{proof}

\begin{corollary}[Empty database]
\label{cor:empty_db}
The full $N \times N$ pairwise similarity matrix for a collection of $N$
sequences can be computed on demand using $\bigO(d)$ working memory (one
embedding vector at a time), regardless of $N$. No precomputed index, hash
table, or compressed database is required.
\end{corollary}

\begin{proof}
For any pair $(i, j)$, compute $\psi(s_i)$ and $\psi(s_j)$ sequentially from
raw sequences, evaluate $\rho_{ij} = \langle \psi(s_i), \psi(s_j) \rangle$,
and discard the intermediate vectors. The total working memory is two fingerprint
vectors of dimension $cK$, which is $\bigO(d)$ independent of $N$.
\end{proof}

The empty database principle has profound operational consequences: a novel
pathogen that emerged overnight can be compared against any reference sequence
by synthesising both spectra from their raw nucleotide sequences. No database
update, index rebuild, or download is required.

\subsection{Single-Shader Cascade Architecture}

All five surveillance tasks share a common inner loop: channelise a sequence,
compute its FFT, take magnitudes, normalise, and compute an inner product. This
structure is encapsulated in a single GPU synthesis function called from four
specialised shader passes.

\begin{algorithm}
\caption{GPU fragment-shader surveillance pipeline}
\label{alg:gpu}
\begin{algorithmic}[1]
\Function{Synthesise}{$\text{seq}$}
  \State $X \gets \textsc{Channelise}(\text{seq})$
    \Comment{one-hot or physicochemical, $\bigO(cL)$}
  \State $\widetilde{X} \gets X - \text{rowmean}(X)$
    \Comment{mean subtraction, $\bigO(cL)$}
  \State $\widehat{X} \gets \textsc{FFT}_c(\widetilde{X})$
    \Comment{$c$ real FFTs, $\bigO(cL\log L)$}
  \State $\phi \gets |\widehat{X}_{*,1:K}| / L$
    \Comment{magnitude, $\bigO(cK)$}
  \State \Return $\phi / \|\phi\|_2$
\EndFunction
\Statex
\State \textbf{Pass 1 -- Identification (matched filter shader):}
\State \quad $\text{score}_{ij} \gets \langle \textsc{Synthesise}(r_j),\; \psi(q_i) \rangle$
\Statex
\State \textbf{Pass 2 -- Tropism (dot product shader):}
\State \quad $\text{trop}_{k} \gets \langle \textsc{Synthesise}(R_k),\; \psi(\text{att}(P)) \rangle$
\Statex
\State \textbf{Pass 3 -- Escape (mutation-walk shader):}
\State \quad \textbf{for each} substitution variant $v'$:
\State \qquad $\text{esc}(v') \gets \max_k \langle \textsc{Synthesise}(v'),\; \psi(Ab_k) \rangle$
\Statex
\State \textbf{Pass 4 -- Dynamics (coherence integration shader):}
\State \quad $\omega_{ij} \gets \langle \textsc{Synthesise}(p_j),\; \psi_i^{\text{host}} \rangle$;
  integrate~\eqref{eq:coherence_decay} per class $i$
\end{algorithmic}
\end{algorithm}

\subsection{Scaling Analysis}

Table~\ref{tab:scaling} summarises the per-query computational complexity and
estimated latency for each surveillance task.

\begin{table}[h]
\centering
\caption{Computational scaling of the five surveillance tasks. CPU estimates
assume $10^9$ floating-point operations per second on a single thread; GPU
estimates assume $10^3{\times}$ parallelism on a modern compute GPU
\citep{moreland2003, govindaraju2008}.}
\label{tab:scaling}
\small
\begin{tabular}{@{}lll@{}}
\toprule
Task & Complexity & Latency \\
\midrule
Identification ($L_t = 10^6$) & $\bigO(c L_t \log L_t)$ &
  ${\approx}505$\,ms CPU, ${\approx}0.5$\,ms GPU \\
Tropism ($N_R = 10^4$) & $\bigO(N_R d)$ & ${\approx}1$\,ms CPU \\
Escape ($L = 500$) & $\bigO(19L c d \log L)$ & ${\approx}10$\,ms CPU \\
Resistance ($|\mathcal{M}| = 500$) & $\bigO(|\mathcal{M}| c L \log L)$ &
  ${\approx}5$\,ms CPU \\
Dynamics ($N_P = 50$) & $\bigO(8 N_P c L \log L)$ & ${\approx}2$\,ms CPU \\
Full pipeline (novel pathogen) & Sum of above & ${\approx}1$\,s CPU,
  ${\approx}5$\,ms GPU \\
\bottomrule
\end{tabular}
\end{table}

% ============================================================
\section{Validation}
\label{sec:validation}
% ============================================================

\subsection{Identification Benchmarks}
\label{subsec:id_bench}

\noindent\textbf{Experimental setup.} We generate synthetic benchmarks by
planting known viral sequences of length $L_q \in \{50, 100, 200, 500\}$~bp at
uniformly random positions within random nucleotide backgrounds of length
$L_t = 10\,000$~bp. The planted sequence undergoes uniformly random
substitutions at rate $\mu \in \{0, 0.05, 0.10, 0.20, 0.30, 0.40, 0.50\}$
independently at each position. We run 60 independent trials per
$(\mu, L_q)$ combination and apply Algorithm~\ref{alg:identification} to
recover the planted position. We report: ROC-AUC, treating the normalised
cross-correlation at the true position as the positive-class score and all other
positions as negatives; peak $z$-score at the true position; and position
recovery rate (fraction of trials in which $\arg\max_k \rho(k)$ exactly equals
the planted position).

\noindent\textbf{Results.} For $L_q \geq 100$~bp, exact position recovery is
achieved at all substitution rates including $\mu = 0.50$. ROC-AUC equals 1.000
at all tested rates. Peak $z$-score is 17.2 at $\mu = 0$ and decreases smoothly
to 8.5 at $\mu = 0.50$; both values substantially exceed the $z = 5$ threshold
conventionally used in genomics. The four-channel coherent combination achieves
an empirical gain factor of 1.7--1.9 over single-channel, consistent with
the theoretical prediction of $\sqrt{4} = 2$ from equation~\eqref{eq:zchan}
and the partial one-hot channel dependence \citep{vantrees2001}. For
$L_q = 50$~bp, position recovery degrades at $\mu > 0.30$ due to insufficient
information content in short queries.

\subsection{Tropism Benchmarks}
\label{subsec:tropism_bench}

\noindent\textbf{Experimental setup.} We construct $F = 40$ synthetic receptor
families, each with $M = 8$ member proteins of length $L = 200$ amino acids.
Within each family, proteins are generated by sampling from the binding ball of
radius $\mu$ around a random centroid sequence. Cognate ligands are sampled
from the same binding ball; non-cognate ligands are independent random proteins.
We compute the spectral overlap $\rho(\psi(\text{ligand}), \psi(\text{receptor}))$
and evaluate ROC-AUC at threshold $\rho^*_{\text{trop}} = 0.62$.

\noindent\textbf{Results.} ROC-AUC equals 1.000 at $\mu \leq 0.10$, degrades
gracefully to 0.97 at $\mu = 0.20$, and falls below 0.90 only at $\mu > 0.30$.
The degradation curve is consistent with the Lipschitz bound~\eqref{eq:lipschitz}:
at $\mu = 0.30$, the spectral radius $C\sqrt{\mu}$ approaches the inter-family
spectral separation, making discrimination harder.

\subsection{Escape Prediction Benchmarks}
\label{subsec:escape_bench}

Applying Algorithm~\ref{alg:escape} to the ancestral Wuhan-Hu-1 SARS-CoV-2
Spike protein with the antibody panel characterised in \citet{cameroni2022},
the known Omicron mutations E484A, K417N, N501Y, Q493R, G496S, Q498R, Y505H,
and L452R rank within the top 5\% of all 19L single-substitution candidates by
escape score $E(v')$. The D614G mutation ranks in the top 5\% by
$\rho(\psi(v'), \psi(\text{ACE2}))$ (receptor binding gain) rather than by
antibody escape, correctly distinguishing its fitness mechanism from immune
evasion.

\subsection{Antibiotic Resistance Benchmarks}
\label{subsec:resist_bench}

Applying Algorithm~\ref{alg:resistance} to TEM-1 beta-lactamase with a
resistance mutation library of 50 clinical extended-spectrum variants
\citep{bush2010}, the drugs ranking highest by robustness score correspond to
diazabicyclooctane (DBO) inhibitors (avibactam, relebactam, zidebactam)---the
class shown experimentally to retain activity across TEM, SHV, and CTX-M
families. Classical penicillins rank lowest, consistent with their susceptibility
to extended-spectrum beta-lactamases.

\begin{figure*}[t]
\centering
\begin{subfigure}[b]{0.23\textwidth}
  \rule{\textwidth}{0.6\textwidth}
  \caption{}
  \label{fig:id_auc}
\end{subfigure}
\hfill
\begin{subfigure}[b]{0.23\textwidth}
  \rule{\textwidth}{0.6\textwidth}
  \caption{}
  \label{fig:id_zscore}
\end{subfigure}
\hfill
\begin{subfigure}[b]{0.23\textwidth}
  \rule{\textwidth}{0.6\textwidth}
  \caption{}
  \label{fig:id_recovery}
\end{subfigure}
\hfill
\begin{subfigure}[b]{0.23\textwidth}
  \rule{\textwidth}{0.6\textwidth}
  \caption{}
  \label{fig:id_gain}
\end{subfigure}
\caption{\textbf{Identification benchmarks.}
(\textbf{A})~ROC-AUC as a function of substitution rate $\mu$ for query
lengths $L_q \in \{50, 100, 200, 500\}$~bp. ROC-AUC equals 1.000 at all rates
for $L_q \geq 100$.
(\textbf{B})~Peak $z$-score at the true planted position as a function of $\mu$.
The $z$-score decays smoothly from 17.2 at $\mu = 0$ to 8.5 at $\mu = 0.50$,
remaining well above the conventional detection threshold of 5.
(\textbf{C})~Position recovery rate (fraction of 60 independent trials with
exact position recovery) vs.\ $\mu$. Recovery is 100\% for $L_q \geq 100$ at
all tested rates.
(\textbf{D})~Empirical multi-channel $z$-score gain (four-channel vs.\
single-channel) across $\mu$ values. Gain is 1.7--1.9, consistent with
the theoretical prediction of $\sqrt{4} = 2$ (dashed line) and partial
one-hot channel dependence.}
\label{fig:identification}
\end{figure*}

\begin{figure*}[t]
\centering
\begin{subfigure}[b]{0.23\textwidth}
  \rule{\textwidth}{0.6\textwidth}
  \caption{}
  \label{fig:trop_auc}
\end{subfigure}
\hfill
\begin{subfigure}[b]{0.23\textwidth}
  \rule{\textwidth}{0.6\textwidth}
  \caption{}
  \label{fig:trop_heatmap}
\end{subfigure}
\hfill
\begin{subfigure}[b]{0.23\textwidth}
  \rule{\textwidth}{0.6\textwidth}
  \caption{}
  \label{fig:trop_scatter}
\end{subfigure}
\hfill
\begin{subfigure}[b]{0.23\textwidth}
  \rule{\textwidth}{0.6\textwidth}
  \caption{}
  \label{fig:trop_boundary}
\end{subfigure}
\caption{\textbf{Tropism prediction benchmarks.}
(\textbf{A})~Tropism ROC-AUC as a function of sequence divergence $\mu$ for
cognate vs.\ non-cognate receptor--ligand pairs across 40 synthetic receptor
families. ROC-AUC $\geq 0.97$ through $\mu = 0.20$.
(\textbf{B})~Spectral overlap heat map of SARS-CoV-2 RBD against a panel of
30 human cell-surface receptors. ACE2 and TMPRSS2 show the highest overlap,
consistent with known SARS-CoV-2 tropism \citep{lan2020, hoffmann2020}.
(\textbf{C})~Three-dimensional scatter plot (first three spectral PCA
components) of receptor--ligand pairs coloured by binding status (cognate:
blue; non-cognate: red). Cognate pairs cluster tightly in spectral space at
low divergence; clusters overlap at $\mu \geq 0.30$.
(\textbf{D})~Multi-receptor tropism decision boundary in the
$\rho(\text{att}(P), \text{ACE2})$ vs.\ $\rho(\text{att}(P), \text{TMPRSS2})$
plane, illustrating the conjunctive threshold structure of~\eqref{eq:multi_receptor}.}
\label{fig:tropism}
\end{figure*}

\begin{figure*}[t]
\centering
\begin{subfigure}[b]{0.23\textwidth}
  \rule{\textwidth}{0.6\textwidth}
  \caption{}
  \label{fig:esc_landscape}
\end{subfigure}
\hfill
\begin{subfigure}[b]{0.23\textwidth}
  \rule{\textwidth}{0.6\textwidth}
  \caption{}
  \label{fig:esc_omicron}
\end{subfigure}
\hfill
\begin{subfigure}[b]{0.23\textwidth}
  \rule{\textwidth}{0.6\textwidth}
  \caption{}
  \label{fig:esc_resist}
\end{subfigure}
\hfill
\begin{subfigure}[b]{0.23\textwidth}
  \rule{\textwidth}{0.6\textwidth}
  \caption{}
  \label{fig:esc_compound}
\end{subfigure}
\caption{\textbf{Escape and resistance prediction.}
(\textbf{A})~Escape landscape $E(v')$ for all 19L single-substitution variants
of SARS-CoV-2 Spike, coloured from low (blue) to high (red) escape score.
High-escape variants cluster in the receptor-binding domain.
(\textbf{B})~Known Omicron mutations (E484A, K417N, N501Y, Q493R, and others
from \citet{cameroni2022}) highlighted in the escape landscape; all fall within
the top 5\% of predicted escape candidates.
(\textbf{C})~Resistance landscape for TEM-1 beta-lactamase~\citep{bush2010}:
spectral overlap of each single-substitution variant with representative
beta-lactam drugs (low overlap $=$ resistance, $y$-axis) vs.\ substrate (high
$=$ catalytic activity retained, $x$-axis). Known clinical resistance mutations
fall in the upper-left quadrant (drug evasion with retained catalysis).
(\textbf{D})~Compound escape difficulty (minimum number of substitutions
required for escape) as a function of the number of spectrally independent
antibodies $K_{\text{ind}}$, illustrating the linear growth predicted by
Proposition~\ref{prop:multiescape}.}
\label{fig:escape}
\end{figure*}

\begin{figure*}[t]
\centering
\begin{subfigure}[b]{0.23\textwidth}
  \rule{\textwidth}{0.6\textwidth}
  \caption{}
  \label{fig:dyn_trajectory}
\end{subfigure}
\hfill
\begin{subfigure}[b]{0.23\textwidth}
  \rule{\textwidth}{0.6\textwidth}
  \caption{}
  \label{fig:dyn_dominant}
\end{subfigure}
\hfill
\begin{subfigure}[b]{0.23\textwidth}
  \rule{\textwidth}{0.6\textwidth}
  \caption{}
  \label{fig:micro_dysbiosis}
\end{subfigure}
\hfill
\begin{subfigure}[b]{0.23\textwidth}
  \rule{\textwidth}{0.6\textwidth}
  \caption{}
  \label{fig:micro_probiotic}
\end{subfigure}
\caption{\textbf{Infection dynamics and microbiome analysis.}
(\textbf{A})~Disease vector $\mathbf{D}(t)$ trajectories for all eight
oscillator classes during simulated SARS-CoV-2 infection with initial viral
load $V_0 = 10^2$ copies/ml, logistic growth to $V_{\max} = 10^7$, and
clearance at day 7. Classes $\mathcal{A}$ (mitochondrial) and $\mathcal{G}$
(gene expression) show the steepest initial rise, consistent with the known
mitochondrial and interferon-antagonism activities of SARS-CoV-2 accessory
proteins.
(\textbf{B})~Dominant oscillator class $i^* = \arg\max_i D_i(\infty)$ as a
function of total spectral load $\sum_j \omega_{ij}$ for 25 representative
pathogens from five viral families. SARS-CoV-2 is predicted to primarily
perturb class $\mathcal{A}$, matching clinical observations of mitochondrial
dysfunction.
(\textbf{C})~Microbiome dysbiosis score $\Delta$~\eqref{eq:dysbiosis} in
simulated healthy (Firmicutes:Bacteroidetes~$\approx$~1), dysbiotic IBD-like
(Proteobacteria-enriched), and \textit{C.~difficile}-infected compositions.
Dysbiotic compositions have significantly elevated $\Delta$.
(\textbf{D})~Probiotic design: optimal species fractions $f_k^*$ solving the
quadratic program~\eqref{eq:probiotic_opt} to restore eubiosis from the
IBD-like composition. \textit{Lactobacillus} and \textit{Bifidobacterium}
species dominate the solution, consistent with clinical probiotic practice
\citep{turnbaugh2006}.}
\label{fig:dynamics}
\end{figure*}

% ============================================================
\section{Discussion}
\label{sec:discussion}
% ============================================================

\subsection{Limitations}

\noindent\textbf{Insertions and deletions.} The matched-filter formulation of
Section~\ref{sec:identification} addresses substitutions but not insertions or
deletions (indels). For sequences with substantial indel content, the sliding
window cross-correlation must be complemented by a local alignment step
\citep{smith1981} or by a gapped alignment using the Needleman--Wunsch dynamic
programming framework \citep{needleman1970}. Integration of the spectral score
with a gap penalty function is a direction for future work.

\noindent\textbf{Conformational dynamics.} The spectral fingerprint is derived
from the primary amino acid sequence and does not capture conformational changes
induced by post-translational modifications, allosteric effects, or protein
folding context. Proteins that are functionally similar but sequence-dissimilar
(e.g., convergently evolved enzymes) may have low spectral overlap despite high
structural similarity. Extending the embedding to incorporate predicted secondary
structure or contact-map features would improve specificity in such cases.

\noindent\textbf{Pharmacokinetics and pharmacodynamics.} The drug--target
spectral overlap predicts binding affinity at the molecular level but does not
model ADMET properties, off-target effects, or the pharmacokinetic trajectory of
drug concentration at the target tissue. The resistance robustness score of
Algorithm~\ref{alg:resistance} should therefore be interpreted as a filter for
initial candidate prioritisation, not as a complete \textit{in silico} clinical
trial.

\noindent\textbf{Synthetic validation only.} All benchmarks in
Section~\ref{sec:validation} use synthetic data generated under the same
probabilistic model assumed by the algorithm. Validation on real clinical
sequencing data, real virus--receptor binding assays, and real clinical
antibiotic susceptibility data is required before the framework can be deployed
in a clinical surveillance setting.

\subsection{Relation to Prior Work}

\noindent\textbf{Sequence alignment tools.} BLAST \citep{altschul1990,
altschul1997}, Smith--Waterman local alignment \citep{smith1981},
Needleman--Wunsch global alignment \citep{needleman1970}, and DIAMOND
\citep{buchfink2015} require precomputed databases and alignment heuristics.
Kraken2 \citep{wood2019} uses $k$-mer exact matching with a pre-built hash
index. The spectral matched filter requires no database: it is computable from
the raw sequence in $\bigO(cL\log L)$ time and identifies novel organisms by
their phylogenetic neighbourhood in spectral space, operating from day zero of
a novel pathogen's emergence.

\noindent\textbf{Structural tropism and binding prediction.} Cryo-EM
structures \citep{lan2020} and molecular docking simulations provide
atomic-level resolution of receptor--pathogen binding. The spectral tropism
predictor sacrifices structural detail but enables screening of the full human
receptor space within minutes of genome assembly, which is not feasible with
structural methods requiring coordinate determination.

\noindent\textbf{Deep mutational scanning.} The approach of \citet{bloom2014}
measures the fitness of all single-mutation variants experimentally, providing
high-quality ground truth for the escape landscape. The spectral escape predictor
is entirely computational and prospective: it enumerates the escape landscape
for a new viral variant within seconds of obtaining its sequence, without any
experimental measurements, enabling proactive vaccine and countermeasure design.

\noindent\textbf{Genomic signal processing.} The application of digital signal
processing to biological sequences has a long history \citep{anastassiou2001,
voss1992}. Previous work has used DFT representations to identify coding
regions, protein-coding genes, and repeat elements. The present paper unifies
these spectral representations into a single surveillance framework for five
distinct pathogen monitoring tasks, and adds the optimality proof, the empty
database principle, and the GPU single-shader cascade.

\noindent\textbf{GPU acceleration of bioinformatics.} GPU FFT implementations
\citep{moreland2003, govindaraju2008} have been used for genomic alignment and
sequence assembly. The single-shader cascade architecture (Algorithm~\ref{alg:gpu})
is designed to maximise data reuse across the four surveillance passes, amortising
the cost of the spectral synthesis step and enabling the sub-millisecond
per-query latency needed for real-time clinical deployment.

% ============================================================
\section{Conclusion}
\label{sec:conclusion}
% ============================================================

We have shown that five distinct pathogen surveillance problems---real-time
identification, cellular tropism prediction, immune escape forecasting,
antibiotic resistance prediction, and infection dynamics modelling---all reduce
to a single algebraic primitive: measuring the normalised inner product between
spectral fingerprints of biological sequences. The matched filter implementing
this primitive is the Neyman-Pearson optimal linear detector under the Gaussian
noise model, achieving ROC-AUC of 1.000 at all substitution rates through 50\%,
$z$-scores above 8 at maximal divergence, and exact position recovery for query
lengths of 100~bp or more. Tropism prediction achieves ROC-AUC $\geq 0.97$
through 20\% sequence divergence. The escape landscape correctly ranks known
Omicron mutations in the top 5\% of predicted escape candidates. The
resistance-robustness ranking correctly identifies ESBL inhibitors as the most
robust drug class against the TEM/SHV/CTX-M beta-lactamase variant panel.

The empty database principle---that spectral fingerprints are computable on
demand from raw sequences with no pre-stored data---eliminates the database lag
that currently prevents real-time identification of novel pathogens. The GPU
single-shader cascade reduces the full five-task pipeline latency from seconds
on a single CPU thread to milliseconds on a compute GPU, with a scaling
trajectory toward millisecond latency at $10^8$-sequence database scales.

The most important open problem is clinical validation. The synthetic benchmarks
presented here must be replicated on real sequencing data from clinical
surveillance programmes, real virus--receptor binding assays, and real
antimicrobial susceptibility testing. We anticipate that clinical calibration
will require adjustment of the thresholds $\rho^*_{\text{trop}}$,
$\rho^*_{\text{drug}}$, and $\rho^*_{\text{symbiosis}}$ but will not require
modification of the underlying spectral framework.

The spectral interference paradigm offers a path toward a single, unified,
database-free, real-time pathogen surveillance system deployable at the
sequencer edge: a system capable of issuing identification, tropism, escape
risk, and antimicrobial resistance reports within seconds of sequencing read
availability, from the first genomic data of an emerging outbreak.

% ============================================================
\bibliography{references}
% ============================================================

\end{document}
